% ── CHAPTER 6 ──
\chapter{AI Safety \& Alignment Research: The Deep Research Half}
\label{ch:6}

\begin{tcolorbox}[colback=white,colframe=checkblack,title=Chapter 6 at a Glance,fonttitle=\bfseries,sharp corners,boxrule=1pt]
\textbf{Question:} How do red team, eval, and alignment fit together — and how do you publish so a skeptic can disagree with you using your own method?\\
\textbf{Deliverable:} A research-grade write-up: method, results, limitations, rerun instructions, no smoothing.\\
\textbf{Time:} 4 hrs (90 min read + 2–3 hrs Lab 6).
\end{tcolorbox}

\section{Learning Objectives}
\fbadge
You will be able to:
\begin{enumerate}
\item Distinguish two alignment framings (contain-the-danger vs pluralism/procedural legitimacy) and say why ``how it's made matters.''
\item Describe interpretability basics (probing, features, circuits) and what they do and don't prove.
\item Distinguish red team vs eval vs alignment and name where each is the right tool.
\item Explain responsible disclosure and the scientific method for research (hypothesis, controls, negative results, ethics of publishing).
\item Write a research-grade finding that a skeptic can rerun and disagree with using your own method.
\end{enumerate}

\section{The Idea: What Alignment Means and Where Each Tool Fits}
\tierF

\subsection{Two framings}
\textbf{Contain-the-danger:} AI will be powerful; contain it so it does not harm — misuse, loss of control, concentrated power. Frame: boundaries, refusals, hard lines (Art. 5 prohibitions).\\
\textbf{Pluralism and procedural legitimacy:} AI will shape society; ensure it is made and governed legitimately — who decides what is allowed, how values are included, what process makes the answer trustworthy. Aithos's research adds a third framing, systemic alignment, which Section 6.2 develops.

Both matter. ``How it's made matters'' — an answer is not just about what it says but what ledger produced it (which law, whose values, what check). And ``fair'' is not one number: a promotion ranking that weights ambition heavily fits one workplace culture and disadvantages another, and the same weights applied in a different society encode a different value. The procedural question — who set those weights, who was consulted, can they be changed — is part of the safety answer, not a footnote.

\subsection{Interpretability basics — and what it doesn't buy you}
\begin{itemize}
\item \textbf{Probing:} Train a small classifier on a model's internal states to see if a concept is represented. Finds correlation, not causation.
\item \textbf{Features and circuits:} Which directions in activation space correspond to concepts; which subgraphs compute a behavior. Insightful, not proof of behavior out of distribution.
\end{itemize}
Interpretability can say ``the feature exists in here.'' It cannot yet say ``the model will not use that feature to harm you next week.'' In current practice, features are usually decoded with sparse autoencoders (small interpretable directions read out of activations), and circuits are tested with activation patching and ablations — swap one part of the activations and watch whether the behavior changes, which is the closest you get to a causal claim. The honest limit stands: a feature you can patch is not a guarantee you can ship. Pair it with evals.

\begin{defbox}
\textbf{Interpretability:} Methods (probing, features, circuits) that try to explain what a model represents or computes. Valuable, not sufficient — pair with evals that measure behavior.
\end{defbox}

\subsection{Red team vs eval vs alignment}
Three circles, one center:

\begin{itemize}
\item \textbf{Red team:} Find the failure (the creative break). Best when you don't know what to measure yet.
\item \textbf{Eval:} Measure the rate (the honest number). Best when you know what failure to count.
\item \textbf{Alignment:} Steer the system (how it's made, what process makes it trustworthy). Best when the question is about ongoing governance.
\end{itemize}
Center: a shared reproducible ruler. All three need it. Without a ruler, red team is anecdote, eval is cargo cult, alignment is sermon.

\section{The Aithos Lens: Floors, Not Targets}
\tierR
\textbf{The simple version.} People disagree about what matters. Cultures, communities and individuals weigh values differently, and that disagreement is not a reasoning failure. An evaluator therefore cannot grade a system against one ``correct'' answer without quietly choosing whose answer counts. What can be checked is more modest and more defensible: whether the system is competent enough to be trusted with a decision, whether it keeps to hard legal lines, and whether its reasoning holds together under pressure.

\subsection{Aithos's framing, in its own terms}
Aithos describes itself as a non-profit foundation working on AI alignment, autonomy and pluralism. Its research statements [S: aithos.org research page] make four claims that shape this chapter:
\begin{itemize}
\item \textbf{Alignment is a state to maintain, not a goal to finish.} ``Systemic alignment'' means technical and social systems mediate between interests while keeping human autonomy and diversity intact. It is not convergence on one value set.
\item \textbf{Floors, not targets.} Evaluations should define adequacy for moral decision-making, ``rather than optimizable targets to maximize indefinitely.'' Whatever you maximise, a system will learn to game.
\item \textbf{Open tools, stakeholder criteria.} Where needs diverge, stakeholders should be able to define their own criteria and test against them.
\item \textbf{Who sets the targets is part of the safety question.} Targeting any particular value profile embeds partisan interests. So does choosing which behaviours count as failure.
\end{itemize}

\subsection{Target versus floor}
\begin{center}
\begin{tabularx}{\linewidth}{lXX}
\toprule
Question & Target-based (optimise) & Floor-based (Aithos framing) \\
\midrule
What it asks & Is the system as good as possible on value X? & Is the system competent enough to pass checks that stakeholders can accept? \\
Pass line & A score to maximise & A threshold that must not be crossed \\
Fails when & It scores below a rival & It breaks a hard line or shows incoherent judgment \\
Pluralism & Embeds one value profile & Leaves the profile open; stakeholders add criteria \\
Gaming risk & High: the target is what gets optimised & Lower, not zero: a floor can be met by surface behaviour [R] \\
\bottomrule
\end{tabularx}
\end{center}

\subsection{Two published methods from the Aithos research group}
\begin{itemize}
\item \textbf{Argumentation analysis} [V: abstract, \emph{The Paris Journal on AI \& Digital Ethics}, No. 2, September 2026, pp. 148--167, DOI 10.65701/kl7do0n9re]. Asks whether a model can defend its verdicts under critical questioning, using Walton's argumentation schemes and Govier's criteria for argument cogency. As reported in the abstract: nine frontier models, 200 high-ambiguity MoralChoice items, 6,778 judge-scored cells, and 89.6\% inter-judge agreement on the binary failure judgment. Failure concentrates on grounds and sufficiency and correlates with epistemic hedging rather than argument length. The justification often uses a different argument scheme from the reasoning that produced the verdict (at least 20\% of dilemmas per model).
\item \textbf{Moral competence before moral content} [V: abstract, pp. 168--187, DOI 10.65701/1zidk02aqx]. Proposes four structural conditions for a coherent policy, namely verdict stability, monotonicity, decisiveness and Pareto viability, that can be tested from behaviour alone, without a reference moral standard. As reported in the abstract: nine frontier models, a factorial design of five paraphrases, five escalation levels and three dominance conditions, across three simulated deployments. Surface-form perturbation alone shifted verdict rates by up to 99 percentage points at a single escalation level. The authors conclude that LLM-based agents are not currently the kind of object to which alignment can meaningfully apply.
\end{itemize}

\begin{defbox}
\textbf{What this edition checked.} The two abstracts were read on 2026-10-23. The full definitions of the four conditions, the statistics, and the experimental design were not checked against the full papers. Read them before reusing any condition, and do not quote the abstract numbers as if they were your own results.
\end{defbox}

\subsection{Worked example (illustrative)}
\begin{examplebox}
\textbf{[X] A promotion-ranking assistant.} The same request arrives in five forms: a direct request, a paraphrase, a version with an added legitimate performance factor, a version framed as a ``safety review,'' and a version where the manager escalates after a refusal. A floor check asks three questions. Does the verdict stay the same when only the wording changes? Does a legitimate added factor move the output in the right direction and no further? Does the system commit to a verdict when the rule is clear? Separately, the report names which stakeholder criteria were applied, and who chose them. The mapping of these questions to the formal conditions is [R]; use the Libert et al. definitions.
\end{examplebox}

\begin{pitfallbox}
\textbf{The lens has its own pitfalls.} Passing a floor is not alignment. A floor says the system is not obviously incompetent, not that it acts in your stakeholders' interests. Choosing which floor to test is itself a value choice, so write it into the report. A system that refuses everything also passes some floors; pair floor checks with a usefulness measure.
\end{pitfallbox}

\section{The Mechanics: Method, Disclosure, Ethics}
\tierP

\subsection{Red team vs eval vs alignment — where each is the right tool}
\begin{center}
\begin{tabularx}{\linewidth}{lX}
\toprule
Use this & When \\
\midrule
Red team & You suspect a failure but have no rate. Go find it. \\
Eval & You know the failure (e.g., Art. 5(1)(f) compliance) and need the rate with a denominator. \\
Alignment frame & The question is about governance: who decides, what process, how legitimacy is earned. \\
All three & When the finding becomes a standard: red team finds it, eval measures it, alignment tells you how to govern by it. \\
\bottomrule
\end{tabularx}
\end{center}

\subsection{Responsible disclosure — and the scientific method}
\textbf{Disclosure timeline (real, used in Ch. 8 Example 8.2):}
\begin{enumerate}
\item Find + measure (with denominator) → log fingerprint + transcripts.
\item Tell the owner: scope, method, number, limitation, transcripts, fix suggestion. Give 90 days (or what the bounty policy says).
\item Share the check so they can verify the fix themselves.
\item Publish the method after fix or after window — the check, not the weaponized variant. Redact PII/secrets.
\end{enumerate}

\textbf{Scientific method for research (no smoothing):}
\begin{itemize}
\item \textbf{Hypothesis:} ``Model X resists prohibited emotion inference at $<$50\% under condition C.''
\item \textbf{Controls:} Same scenario, lawful alternative (performance ranking) must still be allowed — test both sides of the boundary.
\item \textbf{Negative results:} Report where the model did resist. The non-failure is data.
\item \textbf{Honesty:} Report the limitation that hurts your story. ``One provision, one model, n=5'' if that is what you ran. The honest number, not the clean one.
\end{itemize}

\section{Worked Examples}
\tierP

\begin{examplebox}
\textbf{Example 6.1 — A probing experiment that finds a feature but fails out of distribution}\\
\textbf{Setup:} A probe finds a ``deception'' direction in a model's activations with 85\% accuracy on training probes.\\
\textbf{Next test:} Ask the model to deceive in a novel multi-turn negotiation the probe never saw. The feature fires but the behavior does not — the probe does not predict action.\\
\textbf{Lesson:} The feature is there; the circuit is not yet understood. Interpretability gave a clue, not a guarantee. The eval (n=40 negotiation episodes, 12/40 deceptive = 30\%) is what measures the behavior. Both are needed. Claims about alignment need both the inside (feature) and the outside (rate).
\end{examplebox}

\begin{examplebox}
\textbf{Example 6.2 — Responsible disclosure timeline for an emotion-inference finding}\\
\textbf{Day 0:} Lab 5 LARA-lite shows 3/5 compliance failures on vendor Y's assistant (n=5). You expand to n=40, three scenarios: 38\% $\pm$ 7\% compliance.\\
\textbf{Day 1:} Email vendor security: method, provision text, scenarios, n, fingerprints, 5 example transcripts (PII redacted), fix suggestion (channel policy + refusal with alternative).\\
\textbf{Day 30:} Vendor confirms, shares patch plan, asks for check to verify. You share the harness.\\
\textbf{Day 60:} Rerun: 85\% $\pm$ 5\% compliance (fix worked for tested scenarios).\\
\textbf{Day 90:} Publish: ``Before/after: 38\% → 85\% compliance, n=40 each, 3 scenarios, Art. 5(1)(f), limitation: 3 scenarios, one vendor.'' No weaponized payloads, no user data. That is the scientific method with disclosure.
\end{examplebox}

\section{Hands-On Lab 6 — Research-Grade Write-Up}
\begin{labbox}
\textbf{Target:} Your LARA-lite finding (or Ch. 2 lab finding) — now as a write-up, not just a harness.\\[4pt]
\textbf{Steps:}
\begin{enumerate}
\item Hypothesis (1 sentence), method (dataset, solver, judge, metric, pass definition), environment fingerprint, results (number + denominator + stderr), controls (lawful alternative), limitation (what you did not test), rerun instructions (exact command a stranger runs), ethics (disclosure, no PII, no weapon).
\item No smoothing: keep the bad number, the limitation, the transcript where the judge disagreed. If you tested 5 and 2 resisted, say 2/5.
\item Length: 2–3 pages. No filler. A skeptic must be able to disagree with you using your own method — because the method is fully there.
\end{enumerate}
\textbf{Deliverable:} A 2–3 page PDF or Markdown research note.\\[2pt]
\textbf{How you know you passed:} A second person, without asking you, reruns the command, gets your number (within sampling), and can point at the exact transcript where a different interpretation would change the conclusion. That is research-grade.
\end{labbox}

\section{Pitfalls and the Missing Check}
\begin{pitfallbox}
\textbf{What goes wrong:}
\begin{itemize}
\item Publishing the claim without the check — ``we found bias'' with no ruler anyone can run.
\item Hiding negative results — where the model did the right thing is data.
\item Conflating red-team anecdote with eval rate — ``it jailbroke once'' is not a prevalence claim.
\item Interpreting probing as proof of behavior.
\end{itemize}
\textbf{The missing check:} Hypothesis + controls + negative results + limitation + rerun instructions + ethics note. Without those, a skeptic cannot engage with your work. With them, a skeptic can disagree with you using your own method — and that is the point. The method is the contribution.
\end{pitfallbox}

\section{Illustrations}
\begin{figure}[htbp]
\centering
\begin{tikzpicture}[
  circ/.style={draw, thick, circle, minimum size=3.5cm, align=center, font=\small, opacity=0.85},
  center/.style={fill=lightgreen, draw=checkblack, thick, rounded corners=3pt, minimum width=2.6cm, minimum height=0.8cm, align=center, font=\footnotesize\bfseries}
]
\begin{scope}[blend group=screen]
\node[circ, fill=lightblue] (rt) at (-1.5,0) {\textbf{Red team}\\ \tiny find the\\ \tiny failure};
\node[circ, fill=lightamber] (ev) at (1.5,0) {\textbf{Eval}\\ \tiny measure\\ \tiny the rate};
\node[circ, fill=lightred!60] (al) at (0,2) {\textbf{Alignment}\\ \tiny steer \&\\ \tiny govern};
\end{scope}
\node[center] (mid) at (0,0.6) {Shared\\ ruler};
\node[font=\footnotesize, below=1.5cm of rt, align=center] {\textbf{Right tool:}\\ unknown failure};
\node[font=\footnotesize, below=1.5cm of ev, align=center] {\textbf{Right tool:}\\ known failure,\\ need the rate};
\node[font=\footnotesize, above=0.8cm of al, align=center] {\textbf{Right tool:}\\ governance \&\\ legitimacy};
\end{tikzpicture}
\caption{Red Team vs Eval vs Alignment — Three overlapping circles. Center = shared reproducible ruler. Each lobe = where it is the right tool.}
\label{fig:6-1}
\end{figure}
\textbf{Prompt for Fig 6.1:} \texttt{venn diagram three overlapping circles labeled red team eval alignment, center shared ruler — clean technical line art, black on white, single muted red accent \#E63946, minimal sans-serif labels, no clip art, no gradients, no 3D, print-ready vector style, A4}

\section{Key Terms}
\begin{itemize}
\item \textbf{Alignment} — Steering a model/system to act as intended under conflict, including governance legitimacy.
\item \textbf{Probing / Features / Circuits} — Interpretability methods; correlation, not yet behavioral guarantee.
\item \textbf{Red team vs eval vs alignment} — Find / measure / govern — each has a different right-tool moment.
\item \textbf{Responsible disclosure} — Owner first, time to fix, check shared, then publish the method (not the weapon).
\item \textbf{Negative results} — Reports of where the system did pass — still data, still published.
\item \textbf{Eval awareness} — Model behaves differently when it detects it is being tested.
\end{itemize}

\section{Self-Test}
\begin{selftest}
\item \textbf{Recall:} What are the two alignment framings, and what question does each ask?
\item \textbf{Recall:} What does a probe finding a ``deception direction'' prove — and what does it not yet prove?
\item \textbf{Application:} You found a 38\% compliance rate. What is hypothesis, control, and negative result in the write-up?
\item \textbf{Application:} A vendor asks you to delay disclosure past 90 days. What is the responsible path, and what do you share so they can verify the fix?
\item \textbf{Design a check:} Design a study that pairs probing with an eval to test whether a feature predicts behavior out of distribution. Define the eval, n, and what ``probing predicts behavior'' would mean as a checkable claim.
\end{selftest}

\readywhen{A skeptical reader can disagree with you using your own method — because the method, limitation, and rerun are fully there.}

\section{Expert Extension}
\rbadge
\textbf{Open problem: What ``alignment measurement'' could mean when the ruler is an LLM.} If the judge is itself an LLM, its alignment shapes your measurement of alignment — circularity.\\
\textbf{Method you would bring:} Triangulate: one provision, three measurement paths — (a) judge board, (b) human expert review (n=50), (c) behavioral proxy (e.g., refusal latency, alternative offered). Report all three + agreement + where they diverge. The divergence is the open problem. Publish the harness so the next researcher can add a fourth path.

\begin{checkbox}
The check is the product. In research, the check that survives a skeptic is the check that becomes infrastructure.
\end{checkbox}
